% Research note --- synthesis or proof work produced from promoted concepts.
% Copy it into 70 Research/20 Research Notes/<thread>/ and replace every <...> placeholder.
\documentclass[a4paper,10pt]{article}
\IfFileExists{preamble.tex}{\input{preamble.tex}}{\IfFileExists{../../../../../60 Templates/preamble.tex}{\input{../../../../../60 Templates/preamble.tex}}{\IfFileExists{../../../../60 Templates/preamble.tex}{\input{../../../../60 Templates/preamble.tex}}{\input{../../../60 Templates/preamble.tex}}}}
\begin{document}
\ArtifactHeader{\ProjectCode}{<Research-map key>}{<what this research note establishes>}{<YYYY-MM-DD>}%
  {<promoted concepts, meeting session records and applicable Source IDs with exact locators>}

\section{Question and status}
<The Research Question or Claim, current status, assumptions and dependencies.>

\section{<The first idea>}

\begin{defn}[<term>]
<The definition in the project's notation, plus any alternative source notation actually used.>
\end{defn}

\begin{example}
<The smallest computation that makes the definition concrete, worked rather than cited.>
\end{example}

\begin{thm}[<name>]
<The statement with its hypotheses spelled out, since later arguments depend on them.>
\end{thm}

\begin{proof}
<The argument and its adopted status. Cite the exact locator when a source supplies the proof.>
\end{proof}

\begin{remark}
<Why a hypothesis is there, what the notation hides, or how this connects back to a one-variable
result already known.>
\end{remark}

\begin{warning}
<The mistake the idea invites, with the counterexample that settles it. Warnings share the
theorem counter, so this one is citable like any other numbered thing.>
\end{warning}

\section{<The second idea>}
<One section per idea, in dependency order or the relevant source order.>

\begin{checkyourself}
  \item <Something to rebuild unaided --- the test of whether this landed.>
  \item <The unresolved step or check that determines the next research pass.>
\end{checkyourself}
\end{document}
